% =====================================================================
%  PREAMBLE — định dạng theo file mẫu "Tham khao mau bao cao bai tap 1.pdf"
%  Lớp scrartcl (KOMA-Script), cỡ chữ 13pt, biên dịch bằng XeLaTeX
% =====================================================================

% ---------- Font và ngôn ngữ ----------
\usepackage{fontspec}
\setmainfont{texgyretermes}[
  Extension      = .otf,
  UprightFont    = *-regular,
  BoldFont       = *-bold,
  ItalicFont     = *-italic,
  BoldItalicFont = *-bolditalic
]
\setmonofont{DejaVuSansMono}[
  Extension      = .ttf,
  BoldFont       = *-Bold,
  ItalicFont     = *-Oblique,
  BoldItalicFont = *-BoldOblique,
  Scale          = 0.85
]

\usepackage{babel}
\babelprovide[import, main]{vietnamese}

% ---------- Trang và khoảng cách ----------
\usepackage[a4paper, left=3cm, right=2cm, top=2cm, bottom=2cm]{geometry}
\usepackage{setspace}
\setstretch{1.3}
\usepackage{indentfirst}
\setlength{\parindent}{1.27cm}
\setlength{\parskip}{3pt plus 1pt}
\emergencystretch=2em
\pagestyle{plain}

% ---------- Tiêu đề mục: in đậm, cùng cỡ chữ thân bài (như mẫu) ----------
\setkomafont{disposition}{\normalfont\bfseries}
\setkomafont{section}{\normalsize}
\setkomafont{subsection}{\normalsize}
\setkomafont{subsubsection}{\normalsize}
\setkomafont{paragraph}{\normalsize}
\renewcommand*{\sectionformat}{\thesection.\enskip}
\renewcommand*{\subsectionformat}{\thesubsection.\enskip}
\RedeclareSectionCommand[beforeskip=8pt, afterskip=4pt]{section}
\RedeclareSectionCommand[beforeskip=6pt, afterskip=3pt]{subsection}

% ---------- Mục lục: dòng mục cấp 1 không đậm, có chấm dẫn ----------
\DeclareTOCStyleEntry[
  linefill         = \TOCLineLeaderFill,
  entryformat      = \normalfont,
  pagenumberformat = \normalfont
]{tocline}{section}

% ---------- Toán, bảng, hình ----------
\usepackage{amsmath, amssymb}
\usepackage{siunitx}
\sisetup{output-decimal-marker = {,}}
\usepackage{multirow}
\usepackage{array}
\usepackage{tabularx}
\renewcommand{\tabularxcolumn}[1]{m{#1}}
\newcolumntype{Y}{>{\centering\arraybackslash}X}
\newcolumntype{L}{>{\raggedright\arraybackslash}X}
\newcolumntype{P}[1]{>{\raggedright\arraybackslash}m{#1}}
\setlength{\arrayrulewidth}{0.4pt}
\renewcommand{\arraystretch}{1.15}
\usepackage{graphicx}
\usepackage{float}
\usepackage{enumitem}
\setlist{itemsep=1pt, topsep=3pt}
\usepackage{xcolor}
\usepackage{tikz}
\usetikzlibrary{calc}

% Chú thích như mẫu: "Bảng 1. ..." trên bảng, "Hình 1. ..." dưới hình, không in đậm
\usepackage{caption}
\captionsetup{labelsep=period, labelfont=normalfont, textfont=normalfont,
  font=normalsize, justification=centering, skip=4pt}
\captionsetup[table]{position=top}
\captionsetup[figure]{position=bottom}

% Ghi chú ngay dưới bảng
\newcommand{\GhiChuBang}[1]{\par\vspace{2pt}{\footnotesize\raggedright #1\par}}

% ---------- Mã QR cho liên kết website / GitHub ----------
\usepackage{qrcode}

% ---------- Liên kết (để cuối) ----------
\usepackage{xurl}  % đường dẫn dài tự xuống dòng
\usepackage[colorlinks=true, linkcolor=black, citecolor=black, urlcolor=blue!55!black]{hyperref}

% \WebURL, \GitHubURL: in liên kết bấm được (định nghĩa link ở config/thong_tin.tex)
\newcommand{\WebURL}{\expandafter\url\expandafter{\LinkWeb}}
\newcommand{\GitHubURL}{\expandafter\url\expandafter{\LinkGitHub}}
% \QRLink{chiều cao}{macro chứa link}
\newcommand{\QRLink}[2]{\edef\qrtmp{\noexpand\qrcode[height=#1]{#2}}\qrtmp}

% ---------- Tên tiếng Việt cho các thành phần ----------
\AtBeginDocument{%
  \renewcommand{\figurename}{Hình}%
  \renewcommand{\tablename}{Bảng}%
  \renewcommand{\contentsname}{MỤC LỤC}%
  \renewcommand{\refname}{Tài liệu tham khảo}%
}

% =====================================================================
%  MACRO
%  \TODO{...}, \dien : giữ chỗ (báo cáo hiện không còn dùng)
%  \HinhAnh[độ rộng]{file trong figures/}{chú thích}{nhãn}
% =====================================================================
\newcommand{\TODO}[1]{\textcolor{red}{[\,#1\,]}}
\newcommand{\dien}{\textcolor{red}{\ldots}}

\newcommand{\HinhAnh}[4][0.95\linewidth]{%
  \begin{figure}[H]
    \centering
    \IfFileExists{figures/#2}{%
      \includegraphics[width=#1]{figures/#2}%
    }{%
      \fbox{\parbox[c][5cm][c]{#1}{\centering\TODO{Chèn hình: figures/#2}}}%
    }%
    \caption{#3}
    \label{#4}
  \end{figure}%
}
